\section{Conclusions}
\label{sec:conclusions}

Measured against a known truth on \nHist{} independent synthetic histories, the 99.9 percent operational-risk capital of the loss distribution approach is unstable at the horizons a bank has. The interquartile range of the splice capital is \spIqrMin{} to \spIqrMax{} times the true quantile, the capital stability test is red for all \redTotal{} rung-by-cell rows, and \infLo{} to \infHi{} percent of the histories at tail index 0.96 conclude an infinite mean that does not exist while \finLo{} to \finHi{} percent at 1.1 conclude a finite one.
A delta-method theory of the single loss approximation explains the dispersion to within \thSdMaxErr{} percent and the false conclusion about the mean to within \thShareMaxErr{} points, and shows that the dispersion falls only as the inverse square root of the number of recorded losses: \yrsMin{} to \yrsMax{} years of history would bring the interquartile range to 0.25 times the truth. What each estimation component adds is measured on the same histories; modelling the threshold and using the data above it are necessary, the expert's tail is worth the accuracy of the expert, and the robust fits and a misspecified family do not help on this process.
The next steps are to run the protocol on a severity that is not exactly Pareto above the threshold, to contaminate the histories so that the robust estimators can be tested on what they are for, to vary the level and the frequency to check the dependence of the dispersion on $L$ that the theory predicts, and to propose and test an estimator whose interval is part of the capital it reports.
